\documentclass[11pt, a4paper]{article}

% bfw-style.tex — gemeinsame Style-Definitionen für alle Handouts/Aufgabenhefte
% Voraussetzung: Kompilation mit xelatex (wegen fontspec)

\usepackage[ngerman]{babel}
\usepackage{geometry}
\geometry{a4paper, top=2.2cm, bottom=2.2cm, left=2.3cm, right=2.3cm}

\usepackage{fontspec}
% Fallback-Kette: Source Sans Pro -> Calibri -> Segoe UI -> Latin Modern Sans
% (Latin Modern ist immer mit MiKTeX vorhanden)
\IfFontExistsTF{Source Sans Pro}{%
  \setmainfont{Source Sans Pro}[Ligatures=TeX]%
}{\IfFontExistsTF{Calibri}{%
  \setmainfont{Calibri}[Ligatures=TeX]%
}{\IfFontExistsTF{Segoe UI}{%
  \setmainfont{Segoe UI}[Ligatures=TeX]%
}{%
  \setmainfont{Latin Modern Sans}[Ligatures=TeX]%
}}}
\IfFontExistsTF{Source Code Pro}{%
  \setmonofont{Source Code Pro}[Scale=0.92]%
}{\IfFontExistsTF{Consolas}{%
  \setmonofont{Consolas}[Scale=0.92]%
}{%
  \setmonofont{Latin Modern Mono}[Scale=0.92]%
}}

\usepackage{xcolor}
% Farbpalette: hell, freundlich, modern
\definecolor{bfwPrimary}{HTML}{0EA5A4}    % Petrol/Türkis - Primär
\definecolor{bfwPrimaryDark}{HTML}{0F766E} % dunkler Petrol für Headings
\definecolor{bfwAccent}{HTML}{F59E0B}     % Amber - Akzent
\definecolor{bfwSuccess}{HTML}{10B981}    % Grün - Erfolg / Lösung
\definecolor{bfwWarn}{HTML}{F97316}       % Orange - Achtung
\definecolor{bfwInk}{HTML}{0F172A}        % fast Schwarz
\definecolor{bfwInkSoft}{HTML}{475569}    % gedämpftes Grau
\definecolor{bfwBgSoft}{HTML}{F8FAFC}     % off-white
\definecolor{bfwBgTeal}{HTML}{ECFEFF}     % helles Türkis
\definecolor{bfwBgAmber}{HTML}{FEF3C7}    % helles Gelb
\definecolor{bfwBgRose}{HTML}{FFE4E6}     % helles Rosé für Eselsbrücken

\usepackage[colorlinks=true, linkcolor=bfwPrimaryDark, urlcolor=bfwPrimaryDark]{hyperref}
\usepackage{enumitem}
\setlist[itemize]{leftmargin=*, itemsep=2pt, topsep=2pt}
\setlist[enumerate]{leftmargin=*, itemsep=2pt, topsep=2pt}

\usepackage[most]{tcolorbox}
\usepackage{booktabs}
\usepackage{tikz}
\usetikzlibrary{decorations.pathmorphing, positioning, shapes.geometric, arrows.meta}

% Merkkästchen — wichtige Definition / Kernpunkt
\newtcolorbox{merksatz}[1][]{%
  enhanced, breakable,
  colback=bfwBgTeal,
  colframe=bfwPrimary,
  boxrule=0.6pt,
  arc=4pt,
  left=10pt, right=10pt, top=8pt, bottom=8pt,
  fonttitle=\bfseries\color{bfwPrimaryDark},
  title={\faLightbulb\ Merksatz},
  #1
}

% Eselsbrücke — Mnemoniks
\newtcolorbox{eselsbruecke}[1][]{%
  enhanced, breakable,
  colback=bfwBgRose,
  colframe=bfwAccent,
  boxrule=0.6pt,
  arc=4pt,
  left=10pt, right=10pt, top=8pt, bottom=8pt,
  fonttitle=\bfseries\color{bfwWarn},
  title={Eselsbrücke},
  #1
}

% Beispiel-Box
\newtcolorbox{beispiel}[1][]{%
  enhanced, breakable,
  colback=bfwBgSoft,
  colframe=bfwInkSoft,
  boxrule=0.4pt,
  arc=3pt,
  left=10pt, right=10pt, top=8pt, bottom=8pt,
  fonttitle=\bfseries\color{bfwInk},
  title={Beispiel},
  #1
}

% Lösungs-Box (für Aufgabenhefte)
\newtcolorbox{loesung}[1][]{%
  enhanced, breakable,
  colback=bfwBgAmber!40,
  colframe=bfwSuccess,
  boxrule=0.6pt,
  arc=3pt,
  left=10pt, right=10pt, top=8pt, bottom=8pt,
  fonttitle=\bfseries\color{bfwSuccess!70!black},
  title={Lösung},
  #1
}

% Glossar-Eintrag
\newcommand{\glossareintrag}[2]{%
  \par\noindent\textbf{\color{bfwPrimaryDark}#1} \enspace #2\par\smallskip
}

% Section-Stil — etwas farbiger als Standard
\usepackage{titlesec}
\titleformat{\section}
  {\Large\bfseries\color{bfwPrimaryDark}}
  {\thesection}{0.6em}{}
\titleformat{\subsection}
  {\large\bfseries\color{bfwPrimary}}
  {\thesubsection}{0.5em}{}
\titleformat{\subsubsection}
  {\normalsize\bfseries\color{bfwInk}}
  {\thesubsubsection}{0.4em}{}

% TikZ-Stil "skizzenhaft" — etwas weicher als Rough.js, aber stabil
\tikzset{
  roughbox/.style = {
    draw=bfwPrimaryDark, thick, fill=bfwBgTeal,
    rounded corners=4pt, inner sep=6pt
  },
  roughaccent/.style = {
    draw=bfwAccent, thick, fill=bfwBgAmber,
    rounded corners=4pt, inner sep=6pt
  },
  roughline/.style = {
    draw=bfwInkSoft, thick, -{Stealth[length=2.5mm]}
  }
}

% Bullet-Symbole (bewusst kein FontAwesome — vermeidet Font-Installations-Probleme)
\newcommand{\faLightbulb}{$\bullet$}
\newcommand{\faBookmark}{$\bullet$}

% Header/Footer
\usepackage{fancyhdr}
\pagestyle{fancy}
\fancyhf{}
\renewcommand{\headrulewidth}{0pt}
\fancyfoot[L]{\small\color{bfwInkSoft}\thepage}
\fancyfoot[R]{\small\color{bfwInkSoft} BFW M-064 Beschaffung im Büromanagement}


% Footer für Lernfeld 1
\fancyfoot[L]{\small\color{bfwInkSoft}\copyright\ Can Siebert\,\textperiodcentered\,2026\,\textperiodcentered\,Lernfeld 1 --- Die eigene Rolle im Betrieb mitgestalten}

\title{\vspace*{-1.5cm}\Huge\bfseries\color{bfwPrimaryDark}Aufgabenheft Session~4\\[0.4em]
  \large\color{bfwInkSoft}\textnormal{Den Betrieb präsentieren \& Wiederholung --- Übungen im IHK-Format}}
\author{\textsf{Lernfeld 1 --- Die eigene Rolle im Betrieb mitgestalten \textperiodcentered{} \small\copyright\ Can Siebert\,\textperiodcentered\,2026}}
\date{Selbstlernmaterial}

\begin{document}
\maketitle
\thispagestyle{empty}

\begin{merksatz}[title={\bfseries So nutzen Sie dieses Aufgabenheft}]
Bearbeiten Sie alle Aufgaben zunächst \emph{ohne} in die Lösungen zu schauen. Schreiben
Sie Ihre Antworten in vollständigen Sätzen --- die IHK-Prüfung verlangt formulierte
Begründungen, nicht bloße Stichworte. Beachten Sie die \emph{Operatoren} (nennen,
erläutern, beurteilen, begründen). Aufgabe~4 und Aufgabe~6 sind praktische Übungen ---
planen Sie hier wirklich eine Präsentation und wiederholen Sie Lernfeld~1. Erst danach
öffnen Sie den Lösungsteil ab Seite~\pageref{loesungen} und vergleichen.
\end{merksatz}

\tableofcontents
\vspace{1em}
\hrule

\section{Aufgaben}

\subsection*{Aufgabe~1 --- Zweck, Ziele und Zielgruppe (12 Punkte)}

\medskip
Sie sollen im Betrieb eine kurze Präsentation halten und beginnen mit der Vorbereitung.

\begin{enumerate}[label=(\alph*)]
  \item \textbf{Erläutern} Sie den Zweck einer Präsentation und \textbf{nennen} Sie drei mögliche Ziele. \hfill (3~P.)
  \item \textbf{Begründen} Sie, warum die Zielgruppen-/Adressatenanalyse am Anfang steht. \hfill (3~P.)
  \item \textbf{Nennen} Sie drei Fragen, die Sie über Ihr Publikum klären sollten. \hfill (3~P.)
  \item \textbf{Beurteilen} Sie die Aussage: \emph{„Man kann jede Präsentation für jedes Publikum gleich halten.”} \hfill (3~P.)
\end{enumerate}

\subsection*{Aufgabe~2 --- Vorbereitung und Aufbau (12 Punkte)}

\medskip
\begin{enumerate}[label=(\alph*)]
  \item \textbf{Nennen} Sie die vier Schritte der Vorbereitung in der richtigen Reihenfolge. \hfill (3~P.)
  \item \textbf{Beschreiben} Sie den Aufbau Einleitung --- Hauptteil --- Schluss mit je einer Aufgabe. \hfill (3~P.)
  \item \textbf{Erläutern} Sie den Begriff „roter Faden” und seine Bedeutung. \hfill (3~P.)
  \item \textbf{Erstellen} Sie für eine 10-minütige Präsentation eine grobe Zeitplanung der drei Teile. \hfill (3~P.)
\end{enumerate}

\subsection*{Aufgabe~3 --- Präsentationsmedien und Foliengestaltung (14 Punkte)}

\medskip
\begin{enumerate}[label=(\alph*)]
  \item \textbf{Nennen} Sie je einen Vorteil und einen Nachteil von Beamer, Flipchart und Pinnwand. \hfill (6~P.)
  \item \textbf{Erklären} Sie die 6\texttimes6-Regel und das KISS-Prinzip. \hfill (4~P.)
  \item \textbf{Formulieren} Sie folgende überladene Folie in kurze Stichpunkte um: \emph{„Unser Unternehmen wurde 1998 gegründet und beschäftigt heute über 250 Mitarbeiterinnen und Mitarbeiter an drei Standorten und produziert Büromöbel.”} \hfill (4~P.)
\end{enumerate}

\subsection*{Aufgabe~4 --- Kurzpräsentation eines Betriebs planen (15 Punkte)}

\medskip
Planen Sie eine 3--5-minütige Präsentation über einen Betrieb Ihrer Wahl.

\begin{enumerate}[label=(\alph*)]
  \item \textbf{Erstellen} Sie einen Unternehmenssteckbrief (Branche, Rechtsform, Produkte/Dienstleistungen, Standort, Leitbild/Ziele). \hfill (5~P.)
  \item \textbf{Entwerfen} Sie stichwortartig drei bis vier Folien mit rotem Faden (Einleitung --- Hauptteil --- Schluss). \hfill (5~P.)
  \item \textbf{Nennen} Sie drei Punkte guter Vortragstechnik und zwei Maßnahmen gegen Lampenfieber. \hfill (5~P.)
\end{enumerate}

\subsection*{Aufgabe~5 --- Vortragstechnik und Feedback (10 Punkte)}

\medskip
\begin{enumerate}[label=(\alph*)]
  \item \textbf{Nennen} Sie vier Merkmale eines überzeugenden Vortrags. \hfill (4~P.)
  \item \textbf{Erläutern} Sie, was Lampenfieber ist und wie man damit umgeht. \hfill (3~P.)
  \item \textbf{Nennen} Sie drei Feedbackregeln für die Bewertung einer Präsentation. \hfill (3~P.)
\end{enumerate}

\subsection*{Aufgabe~6 --- Wiederholung Lernfeld~1 (Sessions 1--3) (15 Punkte)}

\medskip
\begin{enumerate}[label=(\alph*)]
  \item \textbf{Nennen} Sie das Gesetz, das die Berufsausbildung regelt, sowie je eine Pflicht der Auszubildenden und der Ausbildenden. \hfill (5~P.)
  \item \textbf{Erklären} Sie den Unterschied zwischen Bedürfnis, Bedarf und Nachfrage und \textbf{nennen} Sie die beiden Pole des einfachen Wirtschaftskreislaufs. \hfill (5~P.)
  \item \textbf{Nennen} Sie je ein Beispiel für eine Personen- und eine Kapitalgesellschaft und \textbf{erläutern} Sie den Unterschied bei der Haftung. \hfill (5~P.)
\end{enumerate}

\vspace{1em}
\noindent\textbf{Punkte gesamt: 78}

\newpage

\section{Lösungen}\label{loesungen}

\subsection*{Lösung Aufgabe~1}
\begin{loesung}
\textbf{(a)} Zweck: Informationen anschaulich, verständlich und überzeugend an ein
Publikum vermitteln. Ziele z.\,B.: informieren (Wissen vermitteln), überzeugen (für eine
Idee gewinnen), aktivieren (zum Handeln/zu einer Entscheidung anregen).

\textbf{(b)} Weil Inhalt, Sprachniveau, Beispiele, Umfang und Medienwahl vom Publikum
abhängen. Nur wer seine Adressaten kennt, kann zielgerichtet und verständlich vortragen.

\textbf{(c)} Fragen z.\,B.: Wer sitzt im Publikum (Vorwissen, Funktion, Alter)? Was
interessiert die Gruppe / welchen Nutzen hat sie? Wie groß ist das Publikum und wie viel
Zeit steht zur Verfügung?

\textbf{(d)} Die Aussage ist falsch. Dieselben Inhalte müssen je nach Publikum
unterschiedlich aufbereitet werden: Vor Auszubildenden mehr Grundlagen und einfache
Beispiele, vor der Geschäftsleitung kürzer, mit Fokus auf Nutzen, Kosten und Zahlen.
\end{loesung}

\subsection*{Lösung Aufgabe~2}
\begin{loesung}
\textbf{(a)} 1.~Thema eingrenzen, 2.~roten Faden entwickeln, 3.~Informationen beschaffen,
4.~Zeit planen.

\textbf{(b)} \emph{Einleitung}: begrüßen, Thema/Ziel nennen, Interesse wecken, Ablauf
ankündigen. \emph{Hauptteil}: Kernaussagen strukturiert und mit Beispielen/Visualisierung
darstellen. \emph{Schluss}: Kernbotschaft zusammenfassen, Fazit/Appell, danken, Fragen
zulassen.

\textbf{(c)} Der rote Faden ist die logische, nachvollziehbare Gliederung, die alle Teile
verbindet. Er sorgt dafür, dass das Publikum folgen kann und die Kernbotschaft erhalten
bleibt.

\textbf{(d)} Beispiel: Einleitung ca.\ 1,5~Min, Hauptteil ca.\ 7~Min, Schluss ca.\ 1,5~Min.
Richtwert: Ein- und Ausstieg zusammen rund 20--30~\% der Zeit.
\end{loesung}

\subsection*{Lösung Aufgabe~3}
\begin{loesung}
\textbf{(a)} \emph{Beamer}: $+$ große Gruppen, Bilder/Animationen; $-$ technikabhängig,
verleitet zu Textwüsten. \emph{Flipchart}: $+$ flexibel, live entwickelbar, kein Strom;
$-$ nur kleine Gruppen, wenig Platz. \emph{Pinnwand}: $+$ Ideen sammeln und ordnen, alle
beteiligen; $-$ vorbereitungsintensiv, wenig Text je Karte.

\textbf{(b)} 6\texttimes6-Regel: höchstens ca.\ sechs Stichpunkte je Folie und ca.\ sechs
Wörter je Stichpunkt. KISS = „Keep it short and simple” --- so einfach und reduziert wie
möglich.

\textbf{(c)} Beispiel: „Gegründet 1998” --- „250+ Mitarbeitende” --- „3 Standorte” ---
„Produziert Büromöbel”. Wenig Text, den Rest erzählt der Vortragende.
\end{loesung}

\subsection*{Lösung Aufgabe~4}
\begin{loesung}
\textbf{(a)} Der Steckbrief muss alle Punkte enthalten: Branche, Rechtsform (korrekt
benannt, z.\,B.\ GmbH, KG, AG --- vgl.\ Session~3), Produkte/Dienstleistungen, Standort(e),
Leitbild/Ziele.

\textbf{(b)} Beispielgliederung: \emph{Einleitung} --- Firmenname und Ziel der
Präsentation; \emph{Hauptteil} --- Steckbrief-Punkte anschaulich (Zahlen als Diagramm,
Bild des Standorts), roter Faden; \emph{Schluss} --- Kernbotschaft, Dank, Fragen. Folien
nach 6\texttimes6-Regel, wenig Text.

\textbf{(c)} Vortragstechnik z.\,B.: frei sprechen mit Stichwortkarten, Blickkontakt zum
Publikum, ruhiges Tempo mit Pausen. Gegen Lampenfieber z.\,B.: gute Vorbereitung/Proben
und tiefes Durchatmen vor Beginn (auch: positive Selbstinstruktion, langsamer Start).
\end{loesung}

\subsection*{Lösung Aufgabe~5}
\begin{loesung}
\textbf{(a)} Merkmale z.\,B.: frei sprechen (statt ablesen), deutliche Stimme und
angemessenes Tempo mit Pausen, Blickkontakt zum Publikum, offene und ruhige Körpersprache.

\textbf{(b)} Lampenfieber ist die Nervosität vor/während eines Auftritts --- normal und
sogar nützlich (macht wach). Umgang: gute Vorbereitung und Proben, tief durchatmen,
langsam beginnen, freundlichen Blickkontaktpunkt suchen, Pausen aushalten.

\textbf{(c)} Feedbackregeln z.\,B.: erst Positives, dann Verbesserungsvorschläge; konkret
und beschreibend bleiben (nicht abwertend); in der Ich-Form sprechen; wertschätzend und
umsetzbar formulieren; wer Feedback erhält, hört zu, ohne sich sofort zu rechtfertigen.
\end{loesung}

\subsection*{Lösung Aufgabe~6}
\begin{loesung}
\textbf{(a)} Die Berufsausbildung regelt das \textbf{Berufsbildungsgesetz (BBiG)}.
Azubi-Pflicht z.\,B.: Lern-/Sorgfaltspflicht, Weisungen befolgen, Berufsschulpflicht.
Ausbildenden-Pflicht z.\,B.: Ausbildungspflicht (Inhalte vermitteln), Ausbildungsmittel
bereitstellen, Vergütung zahlen.

\textbf{(b)} \emph{Bedürfnis} = empfundener Mangel mit Wunsch nach Beseitigung;
\emph{Bedarf} = mit Kaufkraft versehenes Bedürfnis; \emph{Nachfrage} = am Markt konkret
geäußerter Bedarf. Die beiden Pole des einfachen Wirtschaftskreislaufs sind
\textbf{Haushalte} und \textbf{Unternehmen}.

\textbf{(c)} Personengesellschaft z.\,B.\ OHG oder KG; Kapitalgesellschaft z.\,B.\ GmbH
oder AG. Haftung: Bei Personengesellschaften haften Gesellschafter (mindestens teilweise)
persönlich und unbeschränkt mit dem Privatvermögen; bei Kapitalgesellschaften haftet
grundsätzlich nur die Gesellschaft mit ihrem Vermögen (beschränkte Haftung).
\end{loesung}

\end{document}
